\chapter{Reproducibility analysis}
% \setcounter{section}{0}
\label{chapter:reproducibility-analysis}

\textbf{Not clear if this belong in here or a later part}

\textbf{VERY ROUGH, JUST INFORMATION DUMP}

As discussed briefly in \Cref{chapter:intro}, a performance analysis service at Durham may cater to a varied audience. From longstanding HPC users and developers in, traditionally, the physical sciences through to newer communities who are using simulations and machine learning to improve the quantitative methods in their research field. With this in mind, jumping straight into a performance analysis report may not be advised if the reproducibility of the software is highly variable. The perofmrnace analysis service therefore might be an access point to a reproducibility analysis. Once more, this would help integrate the performance analysis service into an ongoing relationship between researchers and ARC. This is a similar to a process included in the \href{https://www.archer2.ac.uk/ecse/calls/}{ARCHER2 eCSE calls}.

It is envisaged that a reproducibility analysis would, at least initially, be a `lighter touch' service than the more rigorous performance analysis. However, a report of each software's reliability could be part of a broader up-skilling of developers at Durham, whilst also signposting developers to RSEs. Once again supporting the primary aims of this service: increasing HPC and software competency within Durham, and facilitating discussions between RSEs, researchers and HPC specialists.

Following discussions with Jeff Young of Georgia Institute of Technology, a `repository audit' could become an initial component of a reproducibility analysis. This repository audit is essentially a check over the practice employed in the Git repository by the developer, and could be seen a bar of entry before using the performance analysis service. The RepoAudit developed by Georgia Tech is hosted \href{https://github.com/orgs/gt-sse-center/repositories}{here} and, whilst it should not be seen as an out right automated `rejection' or `acceptance' to using the performance analysis service, it can be viewed as another tool for the initial checks to a software. The interface and reproducibility analysis component of the service should be seen as a method to clean up the process of the performance analysis, i.e., once the performance analysis begins, the analyst can focus on profiling and tracing.

\textbf{Reproducibility analysis materials from Marion Weinzierl:}
\begin{itemize}
    \item \href{https://www.reprohack.org/}{ReproHack}
    \item \href{https://codecheck.org.uk/}{CodeCheck}
    \item \href{https://sc24.supercomputing.org/program/papers/reproducibility-initiative/}{SC Repro Initiative}
    \item \href{https://www.acm.org/publications/policies/artifact-review-and-badging-current}{ACM Repro Badges}
    \item \href{https://book.the-turing-way.org/reproducible-research/reproducible-research}{Turing Repro Research}
\end{itemize}
Following from these resources the plan is to create a `somewhat' standardised, initial reproducibility analysis methodology for this service.